\section{Justificación e importancia}

Si un \textit{release} mueve el YAML y el \ServiceDomain{} acordado se queda en el modelo, el
desfase se ve en el archivo: un \textit{path} nuevo, un \texttt{operationId} que ya no nombra el
\Behavior{}, un \texttt{\$ref} que apunta a otro \textit{schema}. Farzi (2021) y Bhatia (2022)
describen esa deriva entre producto, integración y gobernanza cuando BIAN y el contrato no se
rastrean. Conviene verla ahí, no cuando el servicio ya está en producción.

Revisar a mano no escala cuando hay decenas de APIs y de socios. El \AlignmentScore{} y las
preguntas P1--P4 (§\ref{sec:planteamiento}) sirven para ordenar esa cola con un número
(Shabbir et al., 2025).

OpenAPI, como artefacto, ya está maduro (Casas et al., 2021). Lo que no tenemos escrito es el
cruce con un \ServiceDomain{} BIAN. El procedimiento y la matriz del Anexo~C son esa pieza.

S y E salen de alineación de esquemas y de similitud semántica. No reemplazan una prueba en
\textit{runtime} ni un sello regulatorio (eje E5, §\ref{sec:estado-arte}). La pregunta del plan es
documental: este contrato, en esta versión, ¿refleja este SD?
